\documentclass[11pt]{article}
\usepackage{amsmath}
\usepackage{amsfonts}
\usepackage{amsthm}
\usepackage[utf8]{inputenc}
\usepackage[margin=0.75in]{geometry}

\title{CSC111 Winter 2026 Project 1}
\author{Catherine Abdul-Samad and Joshua Yeung}
\date{\today}

\begin{document}
\maketitle

\section*{Running the game}
No extra libraries are required beyond the standard Python 3.13.5 library.

\section*{Game Map}
Example game map below:

\begin{verbatim}
 -1 -1  3  4
  7  1  2  5
 -1 -1 -1  6
\end{verbatim}

Starting location is: 1

\section*{Game solution}
List of commands:
\begin{enumerate}
    \item take fob
    \item go west
    \item take wallet
    \item go east
    \item go east
    \item go east
    \item go north
    \item take passport
    \item go south
    \item go south
    \item buy tcard
    \item drop wallet
    \item drop passport
    \item go north
    \item go west
    \item go north
    \item take usb drive
    \item take laptop charger
    \item go east
    \item drop usb drive
    \item drop laptop charger
    \item go south
    \item take lucky mug
    \item go north
    \item drop lucky mug
\end{enumerate}

\section*{Lose condition(s)}
Description of how to lose the game:
The player loses if they run out of moves before returning all required items to the target location.

List of commands (abbreviated):
\begin{verbatim}
["go east", "go west"] * 25
\end{verbatim}

Which parts of your code are involved in this functionality:
\texttt{adventure.py}: \texttt{AdventureGame.check\_lose\_condition} checks if \texttt{self.moves >= self.max\_moves}. Calls to movement commands increment \texttt{self.moves}.

\section*{Inventory}

\begin{enumerate}
\item All location IDs that involve items in the game:
Locations: 1 (Fob), 3 (USB Drive, Laptop Charger), 4 (Passport), 5 (Lucky Mug), 6 (T-Card), 7 (Wallet).

\item Item data:
\begin{enumerate}
    \item For Item 1 (Fob):
    \begin{itemize}
    \item Item name: fob
    \item Item start location ID: 1
    \item Item target location ID: 4
    \end{itemize}
    \item For Item 2 (USB Drive):
    \begin{itemize}
    \item Item name: usb drive
    \item Item start location ID: 3
    \item Item target location ID: 4
    \end{itemize}
    \item For Item 3 (Lucky Mug):
    \begin{itemize}
    \item Item name: lucky mug
    \item Item start location ID: 5
    \item Item target location ID: 4
    \end{itemize}
\end{enumerate}

    \item Exact command(s) that should be used to pick up an item, and the command(s) used to use/drop the item:
    \begin{verbatim}
    take fob
    inventory
    drop fob
    \end{verbatim}

    \item Which parts of your code (file, class, function/method) are involved in handling the \texttt{inventory} command:
    \texttt{adventure.py}: \texttt{AdventureGame.\_handle\_inventory} displays the list. \texttt{AdventureGame.\_grab\_item} adds to inventory. \texttt{AdventureGame.\_handle\_drop} drops items.
\end{enumerate}

\section*{Score}
\begin{enumerate}

    \item Briefly describe the way players can earn score in your game. Include the first location in which they can increase their score, and the exact list of command(s) leading up to the score increase:
    Players earn points by picking up items (5 points per distinct item type) and depositing them in their target location (variable points, e.g., 20 for USB drive). Items with no 
    First location to score: 1 (Bahen 1F).
    Commands: \texttt{take fob}.

    \item Copy the list you assigned to \texttt{scores\_demo} in the \texttt{simulation.py} file into this section of the report:


    \item Which parts of your code (file, class, function/method) are involved in handling the \texttt{score} functionality:
\end{enumerate}

\section*{Enhancements}
\begin{enumerate}
    \item Describe your enhancement \#1 here
    \begin{itemize}
        \item Brief description of what the enhancement is (if it's a puzzle, also describe what steps the player must take to solve it):
        \item Complexity level (choose from low/medium/high):
        \item Reasons you believe this is the complexity level (e.g., mention implementation details, how much code did you have to add/change from the baseline, what challenges did you face, etc.)
        \item Name the parts of the code which are involved in this enhancement
        \item Copy the list you assigned to \texttt{enhancements\_demo} in the \texttt{simulation.py} file into this section of the report:
    \end{itemize}

    % Uncomment below section if you have more enhancements; copy-paste as needed
    %\item Describe your enhancement here
    %\begin{itemize}
    %    \item Basic description of what the enhancement is:
    %    \item Complexity level (low/medium/high):
    %    \item Reasons you believe this is the complexity level (e.g., mention implementation details)
    %    \item Name the parts of the code which are involved in this enhancement
    %    \item Copy the list you assigned to \texttt{enhancements\_demo} in the \texttt{simulation.py} file into this section of the report:
    %\end{itemize}
\end{enumerate}


\end{document}
